\begin{thebibliography}{99}
\bibitem{Kormen}
Маннинг, Рагхаван, Шютце
{\itshape Введение в информационный поиск} --- Издательский дом \enquote{Вильямс}, 2011. Перевод с английского: доктор физ.-мат. наук Д.\,А.\, Клюшина --- 528 с. (ISBN 978-5-8459-1623-4 (рус.))
\bibitem{gost}
\href{http://www.ifap.ru/library/gost/7052008.pdf}{Список использованных источников оформлять нужно по  ГОСТ Р 7.05-2008}
\bibitem{wikipedia}
Русская Википедия // \MYhref{https://ru.wikipedia.org}{https://ru.wikipedia.org}
\bibitem{pyre}
Python Documentation: re --- Regular expression operations // \MYhref{https://docs.python.org/3/library/re.html}{https://docs.python.org/3/library/re.html}
\bibitem{unicode}
Unicode Standard Annex \#29: Unicode Text Segmentation // \MYhref{https://unicode.org/reports/tr29/}{https://unicode.org/reports/tr29/}
\end{thebibliography}
